\documentclass[10pt,a4paper]{article}
\usepackage{amsmath,amssymb,amsthm,graphicx}
% Minimal retained preamble for S3 SIGMA 2022 article 027 cards taken
% from sources/sigma-2022-027-twistor-theory-of-dancing-paths/
% pinned-original-source.tex (doi:10.3842/SIGMA.2022.027, SIGMA 18
% (2022), CC BY 4.0).
%
% Omitted and why: the publisher documentclass{sigma} block (source
% lines 1-7, including the ifpdf probe) and the frontmatter commands -
% the sigma class is unavailable in the toolchain and the assembler
% supplies the article class. The author equation numbering (source
% line 9), the author theorem declarations (source lines 11-19,
% including the starred Theorem* and the definition-style group with
% its brace pair on lines 16 and 19) and the author macro subset used
% by the delivered spans (source lines 41-108: \CP \C \Z \PP \RP \tL
% \tP \R \Rho \p \be \u \z \I \ee \ep \spp \tr \eps, the line-90
% redefinition of \Z to \mathcal{Z}, \U \cI \st \uh \zh) are copied
% verbatim below in source order. Omitted from the macro block: the
% margin-note macros (lines 21-40: \mnotex \mnote \md \tm \hook) and
% unused shorthands (\projb \bilb \bX \mf \bw \L \d \bv \a \b \r \l
% \ov \eqq \bea \eea \hg \hw \hnabla \hk \hrho \hL \hpsi \he \hS \T \E
% \ush \zsh \rk \ms \ent \bs \bn \bnbf \n) - none occurs in the
% delivered spans; the page-style \allowdisplaybreaks (line 111) is
% layout-only. The qed command is already provided by the
% amsthm/assembler baseline; the author's verbatim \renewcommand{\qed}{}
% at the end of one delivered proof compiles unchanged.
\usepackage{amsmath,amssymb,amsthm}
\numberwithin{equation}{section}

\newtheorem{Theorem}{Theorem}[section]
\newtheorem*{Theorem*}{Theorem}
\newtheorem{Corollary}[Theorem]{Corollary}
\newtheorem{Lemma}[Theorem]{Lemma}
\newtheorem{Proposition}[Theorem]{Proposition}
{ \theoremstyle{definition}
\newtheorem{Definition}[Theorem]{Definition}
\newtheorem{Example}[Theorem]{Example}
\newtheorem{Remark}[Theorem]{Remark} }

\newcommand{\CP}{\mathbb{CP}}
\newcommand{\C}{\mathbb{C}}
\newcommand{\Z}{\mathbb{Z}}
\newcommand{\PP}{\mathbb{P}}
\newcommand{\RP}{\mathbb{RP}}
\newcommand{\tL}{\widetilde{L}}
\newcommand{\tP}{\widetilde{P}}
\newcommand{\R}{\mathbb{R}}
\newcommand{\Rho}{\mathrm{P}}
\def\p{\partial}
\def\be{\begin{equation}}
\def\u{\mathfrak{u}}
\def\z{\mathfrak{z}}
\def\I{\mathcal{I}}
\def\ee{\end{equation}}
\def\ep{{\varepsilon}}
\newcommand{\spp}{\mathbb{S}}
\newcommand{\tr}{\operatorname{tr}}
\newcommand{\eps}{\varepsilon}
 \renewcommand{\Z}{\mathcal{Z}}
 \newcommand{\U}{\mathcal{U}}
\newcommand{\cI}{{\mathcal I}}
\newcommand{\st}{\mid }
\newcommand{\uh}{\widehat{u}}
\newcommand{\zh}{\widehat{z}}
\title{Rigidity of the dancing conformal structure (Theorem 2.3)}
\author{Dunajski, Maciej}
\date{doi:10.3842/SIGMA.2022.027; CC BY 4.0}
\begin{document}
\maketitle
\setcounter{section}{2}
\setcounter{Theorem}{2}
\setcounter{equation}{0}
\noindent\textit{Editorial note.} All delivered mathematical text is the author's own, copied from the pinned source file sources/sigma-2022-027-twistor-theory-of-dancing-paths/pinned-original-source.tex; no statement or proof step was rewritten, reordered or removed.\par\medskip
\noindent\textit{Editorial note.} Printed numbers anchored to the published PDF (SIGMA 18 (2022) 027, stream-extracted text sigma22-027.pdf, census pdf\_anchor\_evidence): Definition 2.1/2.2, Theorem 2.3/2.4/2.5 in section 2, Proposition 3.1 in section 3, Lemma 4.1 in section 4, Remark 6.1 in section 6, confirming the sectionwise shared master counter Theorem declared in the author's newtheorem block (source lines 11-19: Corollary, Lemma, Proposition, Definition, Example, Remark all hang on the Theorem counter).\par\medskip
\noindent\textit{Editorial note.} Theorem-family numbers and equation numbers inside the card are locally renumbered (contexts taken from other sections of the source cannot keep their paper numbers under the sectionwise counters); every cross-reference is a symbolic ref and resolves inside the card.\par\medskip
\noindent\textit{Editorial note.} The counted span is the author's named proof environment 'Proof of Theorem \textbackslash{}ref\{theo\_rigid\}' (source lines 449-541), a remote proof declared in section 3 (008 remote-named-proof precedent).\par\medskip
\noindent\textit{Editorial note.} Contexts delivered in source order: double\_fib (2.1), definull (2.2), lagrangian (2.2), the affine\_c statement, the section-3 opening framework prose, the total-null definition, the alpha\_s statement, the beta-surface definition, and the bridge sentence; all locally renumbered.\par\medskip
\noindent\textit{Editorial note.} The proof states (source lines 531-532) that one equality 'follows from a rather complicated computation which we have performed on MAPLE' - the author's own disclosure, kept verbatim.\par\medskip
\noindent\textit{Editorial note.} Context [437,442] ends before the sentence 'The result of Lemma 4.1 fits into the framework of half-flat causal structures' (source lines 443-444), which carries two citations outside this card's cite closure and is not needed by the proof.\par\medskip
\noindent\textit{Editorial scope.} Single environment card cut from the SIGMA article doi:10.3842/SIGMA.2022.027 (sigma-2022-027-twistor-theory-of-dancing-paths); statement and proof are the author's own text at source lines 255-261 and 449-541.\par\medskip
\noindent\textit{Allowed standard prerequisites (editorial).}\begin{itemize}
\item Alpha-surface family (Lemma 4.1). Delivered in context (source lines 413-423).. Load-bearing: the conformal class is anti-self-dual because the alpha-surfaces are null (source lines 450-456).
\item Ward theorem with cosmological constant. Cited in the proof (Ward\_cosmo).. Load-bearing: contact structure on the twistor space implies an Einstein metric (source lines 461-463).
\item Nonlinear graviton theorem. Cited in the proof (Penrose).. Load-bearing: ASD Weyl tensor from the alpha-surface family (source lines 455-456).
\end{itemize}\medskip
\begin{Theorem}
\label{theo_rigid}
If the dancing condition on a path geometry $\U$ defines a conformal structure
on $M=(\U\times \Z)\setminus {\mathcal I}$, then $\U=\RP^2$ with its flat projective
structure, and the conformal structure is represented by the
${\rm SL}(3)$-invariant anti-self-dual Einstein metric on $M={\rm SL}(3)/{\rm GL}(2)$.
\end{Theorem}

\be\label{double_fib}
\U\stackrel{\mu}\longleftarrow
{\mathcal I}\stackrel{\nu}\longrightarrow {\Z},
\ee

\begin{Definition}
\label{definull}
A tangent vector $(\dot{u}, \dot{z})\in T_{(u, z)}M$
is called
{\em null} if the path through $u$ in the direction of $\dot{u}$ contains the turning point of $z$.
\end{Definition}

\be
\label{lagrangian}
{\mathcal L}(u, z, \dot{u}, \dot{z})=0,
\ee

\begin{Proposition}
There exist affine coordinates $\big(x^0, x^1, \zeta_0, \zeta_1\big)$
on $M={\rm SL}(3)/{\rm GL}(2)$ such that two points are dancing if and only if they are null separated
with respect to the anti-self-dual Einsten metric
\be
\label{affine_c}
g={\rm d}\zeta_0{\rm d}x^0+{\rm d}\zeta_1 {\rm d}x^1+(\zeta\cdot {\rm d}x)^2, \qquad \text{where}\quad
\zeta\cdot {\rm d}x\equiv
\zeta_0 {\rm d}x^0+\zeta_1 {\rm d}x^1.
\ee
\end{Proposition}

The incidence relation encoding the double fibration
(\ref{double_fib}) is given by a map $\psi\colon \U\times \Z\rightarrow \R$, where
$\I$ is the inverse image $\psi^{-1}(0)\subset \U\times \Z$.
Let $(u, z)\in M$, i.e., $\psi(u, z)\neq 0$, and let
$(\dot{u}, \dot{z})\in T_{(u, z)} M$. The infinitesimal dancing condition
is equivalent to the existence of a pair $(u^*, z^*)$ such that
\begin{gather*}
%\label{long_dance}
\psi(u, z^*)=0, \quad\ \psi(u^*, z)=0, \quad\ \psi (u^*, z^*)=0, \quad\
\dot{u}\frac{\p\psi}{\p u}(u, z^*)=0, \quad\
\dot{z}\frac{\p\psi}{\p z}(u^*, z)=0,
\end{gather*}
where $z^*$ is the path in the direction $\dot{u}$ containing $u$ and
the turning point $u^*$ of the path $z$. Here the notation $\dot{u}\p/\p u$ denotes a differentiation
along a vector field.

We now
eliminate $z^*$ between the 1st and the 4th condition, and eliminate
$u^*$ between the 2nd and the 5th condition. Substituting the resulting
expressions to the 3rd condition gives one relation between
$(u, z, \dot{u}, \dot{z})$ which is of the form (\ref{lagrangian}).

We shall say that a two-dimensional surface $\sigma\subset M$ is totally null
if all tangent vector fields in $T_m\sigma, m\in M$ are null in the sense of
Definition \ref{definull}.

\begin{Lemma}
\label{alpha_s}
There exists a three-parameter family of totally null surfaces
$($call them $\alpha$-surfaces$)$ in $M$ which correspond to points in
${\mathcal I}$. For any $(l, p)\subset {\mathcal I}$ the corresponding
$\alpha$-surface consists of $(L, P)\in M$ such that
\be
\label{alpha_def}
P\subset l, \qquad p\subset L, \qquad P\neq p, \qquad L\neq l.
\ee
\end{Lemma}

In addition to the $\alpha$ surfaces $M$ also admits two two-parameter families
of null surfaces (call them $\beta$-surfaces) defined as follows.
The first family $\beta_1$ consist of a fixed point
$P\in \U$ (hence two parameters), and all lines $L$ in $\U$ not containing $P$. The second family $\beta_2$ consist of a fixed line $L\subset \U$ and all points $P$ not on $L$.
If we have instead chosen to model $M$ on the points and paths in $(\Z, \z)$,
then the roles of $\beta_1$ and $\beta_2$ would swap, but the notion

We are now ready to prove Theorem \ref{theo_rigid}, and thus establish the uniqueness
of the dancing metric.

\begin{proof}[Proof of Theorem \ref{theo_rigid}]
If there exists a conformal structure on $M$, then, by Lemma \ref{alpha_s}
it is necessarily anti-self-dual (with the right choice of the orientation).
Indeed, there exists a~three-parameter family of null surfaces. All null
surfaces in four dimensions are self-dual (SD) or anti-self-dual (ASD), and we choose an orientation
of $M$ such that the $\alpha$-surfaces are SD. Thus the Weyl
tensor of any metric $g$ in the conformal class $[g]$ is ASD by Penrose's nonlinear graviton theorem
\cite{Penrose} adapted to neutral signature.

Now we shall argue that the conformal structure $[g]$ defined
by the $\alpha$-surfaces of Lemma \ref{alpha_s} must contain an Einstein
metric. To show it, it is enough to demonstrate that the twistor
space~$\I$ of incident pairs ($\alpha$-surfaces) admits a contact structure.
The existence of an Einstein metric will then be a consequence of
the theorem of Ward \cite{Ward_cosmo}. The contact
structure is intrinsically defined by the double-fibration picture
(\ref{double_fib}). The space $\I$ is equipped with two one-dimensional
distributions defining the quotients to $\U$ and $\Z$. The resulting
two-dimensional distribution is non-integrable, and thus it defines
a one-form $\theta$ on $\I$ such that $\theta\wedge {\rm d}\theta\neq 0$.

We have shown (after the proof of Lemma \ref{alpha_s}) that
on top of the SD $\alpha$-surfaces there exist the two, two-parameter
family of $\beta$-surfaces.
Let $(\spp, \epsilon)$, and $(\spp, \epsilon')$ be two symplectic
real vector bundles over $M$ (see \cite{ADM, AHS}). We~shall make use of the two isomorphisms: $TM=\spp\otimes \spp'$, and
${\Lambda^2}_-=\spp\odot \spp$, where~${\Lambda^2}_-$ is the rank-3 vector bundle of
ASD two-forms.
Let the two families of
$\beta$ surfaces correspond to ASD two-forms
$\Sigma_1$ and $\Sigma_2$ in $\Gamma\big({\Lambda^2}_-\big)$, or equivalently to two
sections~$\iota$,~$\tilde{\iota}$ of~$\spp$.

We can rescale $\iota$ so that the corresponding two-form is
closed, and proportional to \mbox{${\rm d}x^1\wedge {\rm d}x^2$} for some functions $\big(x^1, x^2\big)$ on~$M$
which are constant on each $\beta$-surface in the first two parameter family.
The Frobenius integrability conditions now imply the local existence of
two more coordinates such that
$\zeta_0$ and $\zeta_1$ on~$M$ such that $\operatorname{Ker}(\Sigma_1)=
\operatorname{span}\{\partial/\partial \zeta_0, \partial/\partial \zeta_1\}$.
The functions $(\zeta_0, \zeta_1)$ are then the coordinates on the $\beta$-surface.
The corresponding metric takes the form
\be
\label{metric_md}
g={\rm d}\zeta_A\odot {\rm d}x^A+\Theta_{AB}(x, \zeta){\rm d}x^A\odot {\rm d}x^B, \qquad A, B=0, 1
\ee
for some symmetric two-by-two matrix $\Theta$. The anti-self-duality
condition on the Weyl tensor forces the components of $\Theta$ to be at most
cubic in $(\zeta_0, \zeta_1)$, with additional algebraic relations between the components. Imposing the Einstein condition gives
\be
\label{thetaAB}
\Theta_{AB}=\zeta_A\zeta_B+\Rho_{(AB)}-\Gamma_{AB}^C\zeta_C,
\ee
where the coefficients $\Gamma_{AB}^C$ only depend on $\big(x^0, x^1\big)$,
and $\Rho_{AB}$ is the projective Schouten tensor of
affine connection $\nabla$ with connection components~$\Gamma_{AB}^C$.
The $\beta$-distribution defined by~$\Sigma_1$ is parallel, and
\[
\nabla^g \Sigma_1=6{\mathcal A}\otimes \Sigma_1,
\]
where $\nabla^g$ is the Levi-Civita connection of $g$, and the one-form~${\mathcal A}$ is a symplectic connection for the symplectic structure
structure~\cite{DM} given by
\be\label{omega11}
\Omega={\rm d}{\mathcal A}={\rm d}\zeta_A\wedge {\rm d}x^A+\Rho_{AB}{\rm d}x^A\wedge {\rm d}x^B.
\ee
Up to this point our computation has followed \cite{DM}. Now consider the second
family of $\beta$ surfaces given by the ASD two-form $\Sigma_2$, which in our
chosen coordinate system is given by
\[
\Sigma_2= \epsilon^{AB}\big({\rm d}\zeta_A+\Theta_{AC}{\rm d}x^C\big)\wedge \big({\rm d}\zeta_B+\Theta_{BD}{\rm d}x^D\big).
\]
The normalisation condition $\epsilon(\iota, \tilde{\iota})=1$, and
the condition that the distribution defined by $\Sigma_2$ is parallel
together imply that this two-form satisfies
\[
\nabla^g\Sigma_2=-6{\mathcal A}\otimes \Sigma_2,
\]
where ${\mathcal A}$ is given by (\ref{omega11}). We~will need only the skew-symmetric part of
this equation, which gives
\begin{gather*}
0={\rm d}\Sigma_2+6{\mathcal A}\wedge\Sigma_2 =\epsilon^{CD}\nabla_{[A}\Rho_{B]C} {\rm d}x^A\wedge {\rm d}x^B\wedge {\rm d}\zeta_D,
\end{gather*}
where the second equality follows from a rather complicated computation which
we have performed on MAPLE. Therefore this second family of $\beta$-surfaces
exists iff $\nabla_{[A}P_{B]C}=0$, which (for two-dimensional projective structure) is precisely the condition (see, e.g., \cite{BDE}) that the
connection~$\nabla$ is projectively
equivalent to a flat connection, where we can chose
 $\Gamma_{AB}^C=0$ and $\Rho_{AB}=0$. It was shown in~\cite{DM} that
the metric~(\ref{metric_md}) with $\Theta_{AB}$ given by~(\ref{thetaAB})
is invariant under projective changes of connection complemented by an affine
translation of $\xi_A$. Therefore the resulting metric is of the form
(\ref{affine_c}), which completes the proof.
\end{proof}

\begin{thebibliography}{99}
% Card-local bibliography fragment: the \cite keys inside the
% delivered spans of this card. Adapted from the author's
% thebibliography block (source lines 802-896) of
% sources/sigma-2022-027-twistor-theory-of-dancing-paths/
% pinned-original-source.tex, with the \href wrappers rendered as
% plain text (inner content kept) because the card baseline carries
% no hyperref (bibliography typesetting only; no statement or proof
% text is touched).

\bibitem{ADM}
Atiyah M., Dunajski M., Mason L.J., Twistor theory at fifty: from contour
 integrals to twistor strings, \textit{Proc.~A} \textbf{473} (2017), 20170530,
 33~pages, arXiv:1704.07464.


\bibitem{AHS}
Atiyah M.F., Hitchin N.J., Singer I.M., Self-duality in four-dimensional
 {R}iemannian geometry, \textit{Proc.~A} \textbf{362} (1978), 425--461.


\bibitem{BDE}
Bryant R., Dunajski M., Eastwood M., Metrisability of two-dimensional
 projective structures, \textit{J.~Differential Geom.} \textbf{83} (2009),
 465--499, arXiv:0801.0300.


\bibitem{DM}
Dunajski M., Mettler T., Gauge theory on projective surfaces and anti-self-dual
 {E}instein metrics in dimension four, \textit{J.~Geom. Anal.} \textbf{28}
 (2018), 2780--2811, arXiv:1509.04276.


\bibitem{Penrose}
Penrose R., Nonlinear gravitons and curved twistor theory, \textit{Gen.
 Relativity Gravitation} \textbf{7} (1976), 31--52.


\bibitem{Ward_cosmo}
Ward R.S., Self-dual space-times with cosmological constant, \textit{Comm.
 Math. Phys.} \textbf{78} (1980), 1--17.

\end{thebibliography}
\end{document}
